\chapter{Introduction}
\label{chap:intro}

\section{Problem Context}
\label{sec:intro-context}

\par E-commerce has become one of the most important branches of the digital economy, with accelerated growth over the past decade. E-commerce platforms provide consumers with access to millions of products, but this abundance creates a fundamental problem: decision overload. Users spend significant time browsing, comparing, and evaluating products before making a purchase decision, and this process becomes increasingly complex as the supply grows.

\par Recommendation systems have emerged as a natural solution to this problem, employing various algorithmic techniques --- from collaborative filtering and co-occurrence analysis to content-based similarity --- to reduce the search space and present users with relevant products \cite{LLM4Rec2025, ComprehensiveRecSys2024}. These classical systems are computationally efficient and well-studied in the literature, but they present an important limitation: their outputs are opaque. A user who receives a ``similar products'' or ``frequently bought together'' recommendation does not understand why those products are relevant to their specific situation.

\par With the emergence of large language models (LLMs), it has become possible to generate natural language explanations for recommendations, to synthesize other buyers' reviews, and to provide structured comparisons between similar products \cite{LLMEnhancedRecSys2024, ExplainableRecKG2024}. However, integrating LLMs into commercial applications raises significant challenges related to trust, accuracy, and interaction design.

\par A relevant example is Amazon Rufus, Amazon's AI shopping assistant, which serves over 300 million users and has generated estimated incremental sales of \$12 billion \cite{RufusRevenue2025}. Despite these impressive figures, independent studies have shown that Rufus has an accuracy of only 32\% in identifying the most suitable products, a price hallucination rate of 28\%, and an 83\% tendency to recommend Amazon-owned products \cite{RufusWrong2024, RufusConfidentlyWrong2024}. These documented problems underscore the fact that, even at the largest scale, challenges related to factual accuracy, user trust, and transparency remain unsolved.

\par At the consumer perception level, the data is equally concerning: 55\% of consumers state that they do not trust AI chatbot recommendations \cite{ConsumerDistrust2025}, and 42\% perceive e-commerce AI assistants as upselling tools rather than genuine advisors \cite{AIUpsell2025}. Research in human-computer interaction (HCI) suggests that AI interventions timed after a natural user pause are more effective than proactive ones \cite{CHI2024Timing}, which opens the path for reactive interaction models.


\section{Motivation}
\label{sec:intro-motivation}

\par The combination of classical recommendation systems, which excel at fast candidate generation, and LLM models, which can generate explanations and summaries, offers an opportunity to build assistance systems that are simultaneously relevant, explainable, and efficient. However, the practical integration of these components raises essential questions: when should the system intervene? How can it justify recommendations without introducing fabricated information? What architecture provides the best tradeoff between relevance, latency, cost, and user trust?

\par This thesis is motivated by the observation that existing AI shopping assistants tend to be either too intrusive --- intervening automatically without the user's explicit consent --- or insufficiently transparent in how they generate recommendations. This work proposes a \textit{reactive-first} interaction model, in which the system monitors behavioral signals from the user (time spent on a product, scroll depth, interaction with the reviews section) to detect moments of readiness for assistance, but generates AI advice only when the user explicitly requests it.

\par This approach is supported by recent research showing that AI suggestions offered after a natural pause are better received than preemptive ones \cite{CHI2024Timing}, as well as by market data indicating significant consumer resistance to proactive AI chatbots \cite{ConsumerResistanceAI}.


\section{Objectives}
\label{sec:intro-objectives}

\par The main objective of this thesis is the design, implementation, and evaluation of a trustworthy AI shopping assistant that combines classical recommendation methods with explanations generated by large language models in a reactive interaction model. Specifically, the thesis aims to:

\begin{enumerate}
  \item \textbf{Design a hybrid recommendation architecture} that uses classical algorithms (co-occurrence with Jaccard coefficient, content-based similarity, trends with exponential time decay) for fast candidate generation, and an LLM model (Gemini Flash) for re-ranking and generating natural language explanations.

  \item \textbf{Implement a reactive interaction model} based on lightweight behavioral signals --- dwell time, scroll depth, review interaction --- that indicate the user's readiness for assistance without automatically triggering LLM calls.

  \item \textbf{Develop a review summarization mechanism} grounded in real buyer data, which surfaces both majority sentiments and conflicting opinions, addressing the most criticized problem of existing systems such as Rufus.

  \item \textbf{Implement an AI-assisted comparison mode} that allows users to compare 2--3 similar products in a structured format with tradeoff analysis.

  \item \textbf{Evaluate the system} at three levels: recommendation quality (offline metrics), system performance (latency and cost), and user perception (usefulness, trust, intrusiveness) through a structured user study.
\end{enumerate}


\section{Research Questions}
\label{sec:intro-rq}

\par To guide the design and evaluation of the system, this thesis formulates the following research questions:

\begin{description}
  \item[\textbf{RQ1:}] Does a reactive AI assistant produce a higher level of trust and a lower perception of intrusiveness compared to proactive AI suggestions?

  \item[\textbf{RQ2:}] Do LLM summaries grounded in buyer reviews improve the perceived usefulness and credibility of product information?

  \item[\textbf{RQ3:}] Does AI-assisted comparison improve user confidence when choosing between similar products?

  \item[\textbf{RQ4:}] What are the latency and cost tradeoffs of a two-stage recommendation architecture?
\end{description}


\section{Contributions}
\label{sec:intro-contributions}

\par The main contributions of this thesis are:

\begin{enumerate}
  \item A working \textbf{hybrid recommendation architecture} that combines classical algorithms with an LLM explanation layer in a two-stage pipeline (candidate generation under 50ms, LLM re-ranking and explanation under 500ms).

  \item A \textbf{reactive-first interaction model} based on behavioral signals, where the system prepares context and displays subtle indicators, and the LLM call is triggered only at the user's initiative.

  \item A \textbf{review summarization mechanism} grounded in real data, which surfaces conflicts and divergent opinions, not just the majority sentiment.

  \item An \textbf{AI-assisted comparison flow} for similar products within the same category.

  \item A \textbf{complete working prototype} implemented as a full-stack web application with a simulated store, realistic seed data, a recommendation engine, behavioral tracking, and an advisor interface.

  \item A \textbf{systematic evaluation} at three levels: offline recommendation metrics, performance measurements (latency, token cost), and a structured user study.

  \item \textbf{Design guidelines} for building non-intrusive AI assistants in consumer-facing systems.
\end{enumerate}


\section{Thesis Structure}
\label{sec:intro-structure}

\par The thesis is organized into seven chapters, as follows:

\par \textbf{Chapter~\ref{chap:background} --- Background and Related Work} presents the theoretical foundations of classical recommendation systems, the integration of LLMs into recommendation systems, research on trust and explainability in human-AI interaction, and a detailed case study of Amazon Rufus as a reference point for production-scale AI shopping assistants.

\par \textbf{Chapter~\ref{chap:problem} --- Problem Definition and Requirements} formulates the concrete problem addressed by this thesis, defines the functional and non-functional system requirements, establishes the evaluation criteria, and clarifies the project's boundaries and assumptions.

\par \textbf{Chapter~\ref{chap:design} --- System Design} describes the overall system architecture, the two-stage recommendation pipeline, the reactive interaction model, the behavioral tracking mechanism, the review summarization flow and comparison mode, as well as the justification for architectural decisions.

\par \textbf{Chapter~\ref{chap:implementation} --- Implementation} presents the technical details of the prototype: the technology stack, the monorepo structure, the Convex data model, the implemented recommendation algorithms, the LLM integration with Gemini Flash, and the user interface.

\par \textbf{Chapter~\ref{chap:evaluation} --- Evaluation} describes the evaluation methodology and presents results at all three levels: offline recommendation metrics, system performance measurements, and user study results. It includes a discussion of the results, limitations, and threats to validity.

\par \textbf{Chapter~\ref{chap:conclusions} --- Conclusions} summarizes the contributions, answers the research questions based on the obtained results, and proposes future directions for research and development.
